\chapter{Time Series Analysis \& Temporal Anomaly Detection}

\begin{tcolorbox}[colback=primarylight, colframe=primarydark, title={\textbf{What You Will Learn in this Chapter}}]
\begin{itemize}
    \item Foundations of temporal financial data: High-frequency transaction velocity and settlement flows
    \item Classical decomposition: Trend ($T_t$), Seasonality ($S_t$), and Irregular noise ($\epsilon_t$)
    \item Weak stationarity conditions and unit-root testing via Augmented Dickey-Fuller (ADF)
    \item Mathematical structure of Autoregressive Integrated Moving Average (ARIMA) models
    \item Statistical process control: Dynamic rolling-window confidence bounds and CUSUM change-point detection
    \item Hands-on Python implementation of ARIMA residual forecasting and withdrawal burst detection
\end{itemize}
\end{tcolorbox}

\section{The Temporal Dimension of Financial Forensics}

Financial transactions do not occur as independent, identically distributed (i.i.d.) static observations; they are dynamic \index{time series}\textbf{time series} indexed sequentially in time: $\{y_t\}_{t=1}^T$. Forensic investigators analyze temporal series to detect abrupt structural shifts, unusual volume bursts, and money laundering layering spikes.

\begin{definitionbox}{Classical Time Series Decomposition}
\index{Time Series Decomposition}
A financial time series $y_t$ can be decomposed into additive or multiplicative components:
\begin{align}
\text{Additive:} \quad y_t &= T_t + S_t + C_t + \epsilon_t \\
\text{Multiplicative:} \quad y_t &= T_t \times S_t \times C_t \times \epsilon_t
\end{align}
where $T_t$ represents the long-term secular trend, $S_t$ is deterministic periodic seasonality (e.g., end-of-month payroll spikes), $C_t$ is cyclical macroeconomic variation, and $\epsilon_t$ is stochastic white noise.
\end{definitionbox}

\begin{videocard}{IITM BS Lecture: Introduction to Time Series}{https://www.youtube.com/watch?v=odptnlDEz8U}
Understanding temporal dependencies in finance, timestamp sequencing, and frequency aggregation.
\end{videocard}

\begin{videocard}{IITM BS Lecture: Components of Time Series}{https://www.youtube.com/watch?v=0CIdS0syn34}
Decomposing financial signals into underlying trends, seasonal patterns, and residual error terms.
\end{videocard}

\section{Stationarity and the Augmented Dickey-Fuller Test}

A time series is said to be \index{stationarity}\textbf{weakly (covariance) stationary} if its statistical properties do not change over time:
\begin{enumerate}
    \item Constant expected value: $E[y_t] = \mu$ for all $t$.
    \item Constant finite variance: $\text{Var}(y_t) = \sigma^2 < \infty$ for all $t$.
    \item Autocovariance depends only on lag $k$, not on time $t$: $\text{Cov}(y_t, y_{t-k}) = \gamma_k$.
\end{enumerate}

\begin{videocard}{Masterclass Video: Stationarity in Time Series -- Ritvikmath}{https://www.youtube.com/watch?v=oY-j2Wof51c}
Ritvikmath's visual explanation of why non-stationary time series generate spurious correlations and how differencing achieves stationarity.
\end{videocard}

\begin{theorembox}{The Augmented Dickey-Fuller (ADF) Test}
\index{Augmented Dickey-Fuller Test}
To test for the presence of a unit root (non-stationarity), estimate the regression:
\begin{equation}
\Delta y_t = \alpha + \beta t + \gamma y_{t-1} + \sum_{i=1}^p \delta_i \Delta y_{t-i} + \epsilon_t
\end{equation}
The hypothesis test is formulated as:
\begin{itemize}
    \item $H_0: \gamma = 0$ (The series possesses a unit root and is non-stationary).
    \item $H_1: \gamma < 0$ (The series is stationary).
\end{itemize}
If the ADF test-statistic is more negative than the critical value (or $p < 0.05$), we reject $H_0$ and conclude the series is stationary.
\end{theorembox}

\section{ARIMA Modeling for Baseline Forecasting}

To detect temporal anomalies, forensic engines fit an \index{ARIMA}\textbf{Autoregressive Integrated Moving Average (ARIMA$(p, d, q)$)} model on historical baseline behavior.

\begin{formulabox}{ARIMA$(p, d, q)$ Formulation}
For a $d$-times differenced series $y_t' = \nabla^d y_t$:
\begin{equation}
y_t' = c + \sum_{i=1}^p \phi_i y_{t-i}' + \sum_{j=1}^q \theta_j \epsilon_{t-j} + \epsilon_t
\end{equation}
where $\phi_1, \dots, \phi_p$ are autoregressive coefficients, $\theta_1, \dots, \theta_q$ are moving average coefficients, and $\epsilon_t \sim \mathcal{N}(0, \sigma^2)$ is Gaussian white noise.
\end{formulabox}

\begin{videocard}{IITM BS Lecture: Time Series Models}{https://www.youtube.com/watch?v=rHDKwLcsM2o}
Mathematical specification of AR, MA, ARMA, and seasonal ARIMA models for financial forecasting.
\end{videocard}

\begin{videocard}{IITM BS Lecture: Time Series Models -- Hands-On}{https://www.youtube.com/watch?v=3LI8ziBMVJ0}
Practical tutorial using Python statsmodels to fit ARIMA models, analyze ACF/PACF plots, and inspect residuals.
\end{videocard}

\section{Temporal Anomaly Detection: Dynamic Confidence Intervals and CUSUM}

Once an ARIMA model is trained on legitimate baseline history, it generates one-step-ahead forecasts $\hat{y}_t$ with conditional variance $\sigma_t^2$.

\begin{intuitionbox}{The Dynamic Forensic Envelope}
Static thresholds (e.g., ``flag any withdrawal $> \text{Rs.~5,00,000}$'') fail because normal business payroll spikes on the 1st of every month get flagged, while an illicit Rs.~3,00,000 extraction on a quiet Sunday afternoon passes undetected. An ARIMA confidence envelope expands during expected high-volatility hours and contracts during quiet periods, exposing stealthy off-hour transactions.
\end{intuitionbox}

\begin{definitionbox}{CUSUM (Cumulative Sum Control Chart)}
\index{CUSUM}
CUSUM monitors cumulative positive and negative deviations from the baseline mean $\mu_0$:
\begin{align}
S_t^+ &= \max\left(0, S_{t-1}^+ + (y_t - \mu_0) - K\right) \\
S_t^- &= \max\left(0, S_{t-1}^- - (y_t - \mu_0) - K\right)
\end{align}
where $K$ is the allowable tolerance parameter. An anomaly alarm fires if $S_t^+ > H$ or $S_t^- > H$, detecting persistent micro-leakage (such as salami slicing fraud).
\end{definitionbox}

\begin{figure}[htbp]
\centering
\begin{tikzpicture}
\begin{axis}[
    width=0.85\textwidth,
    height=5.2cm,
    xmin=0, xmax=20,
    ymin=0, ymax=10,
    xlabel={Time Step ($t$)},
    ylabel={Cumulative Sum ($S_t^+$)},
    grid=major,
    grid style={dashed, gray!30},
    legend style={at={(0.05,0.95)}, anchor=north west, font=\footnotesize}
]
\addplot[ultra thick, primaryborder] coordinates {
    (0,0) (1,0) (2,0.2) (3,0) (4,0.1) (5,0) (6,0) (7,0.3) 
    (8,0) (9,0.5) (10,1.2) (11,2.1) (12,3.3) (13,4.6) (14,6.0) (15,7.5) (16,9.1)
};
\addplot[dashed, ultra thick, red] coordinates {(0,5.0) (20,5.0)};
\node[fill=red, text=white, rounded corners, font=\tiny\bfseries] at (axis cs:13.5,5.0) {Alarm Threshold ($H=5.0$)};
\legend{$S_t^+$ Cumulative Drift, Alarm Boundary ($H$)}
\end{axis}
\end{tikzpicture}
\caption{CUSUM Control Chart: Detecting Persistent Low-Magnitude Siphoning (Salami Slicing) Crossing the Alarm Limit $H$.}
\end{figure}

\section{Data Science Relevance: Temporal Residual Anomaly Detection}

\begin{lstlisting}[caption={Time Series Residual Anomaly Detection in Python}]
import numpy as np
import pandas as pd
from statsmodels.tsa.arima.model import ARIMA
from statsmodels.tsa.stattools import adfuller

# Simulate 120 days of daily withdrawal volumes (in Rs. Thousands)
np.random.seed(42)
days = pd.date_range(start='2026-01-01', periods=120, freq='D')
base_signal = 100.0 + np.sin(np.arange(120) * (2 * np.pi / 7)) * 15.0 # Weekly cycle
noise = np.random.normal(0, 5, 120)
withdrawals = base_signal + noise

# Inject sudden fraudulent cash drain anomalies on Day 85 and Day 86
withdrawals[85] += 65.0  # Massive spike
withdrawals[86] += 70.0

ts_df = pd.DataFrame({'date': days, 'amount': withdrawals}).set_index('date')

# 1. ADF Test for Stationarity
adf_res = adfuller(ts_df['amount'])
print(f"ADF Statistic: {adf_res[0]:.3f}, p-value: {adf_res[1]:.4e}")

# 2. Fit ARIMA(1, 0, 1) model on baseline (first 75 days)
train_series = ts_df['amount'].iloc[:75]
model = ARIMA(train_series, order=(1, 0, 1))
fitted_model = model.fit()

# 3. Dynamic Residual Forecasting on full series
preds = fitted_model.predict(start=0, end=len(ts_df)-1)
residuals = ts_df['amount'] - preds
residual_std = residuals.iloc[:75].std()

# 4. Flag residuals exceeding 3 standard deviations
ts_df['forecast'] = preds
ts_df['residual'] = residuals
ts_df['anomaly_flag'] = np.abs(residuals) > (3.0 * residual_std)

anomalies = ts_df[ts_df['anomaly_flag']]
print(f"\nDetected {len(anomalies)} Temporal Anomalies:")
print(anomalies[['amount', 'forecast', 'residual']])
\end{lstlisting}

\section*{Supplementary Lecture Videos}
\begin{itemize}
    \item \href{https://www.youtube.com/watch?v=odptnlDEz8U}{W10\_L1: Introduction to Time Series in Financial Forensics}
    \item \href{https://www.youtube.com/watch?v=0CIdS0syn34}{W10\_L2: Components of Time Series (Trend, Seasonality, Noise)}
    \item \href{https://www.youtube.com/watch?v=rHDKwLcsM2o}{W10\_L3: Time Series Models (AR, MA, ARMA, ARIMA)}
    \item \href{https://www.youtube.com/watch?v=3LI8ziBMVJ0}{W10\_L4: Time Series Models -- Hands-On Python Tutorial}
    \item \href{https://www.youtube.com/watch?v=oY-j2Wof51c}{External Masterclass: Stationarity in Time Series (Ritvikmath)}
\end{itemize}

\newpage
\section{Practice Exercises}
\textit{Detailed step-by-step solutions are available in the Appendix.}

\begin{enumerate}
    \item Define covariance stationarity. Given the first-order autoregressive process:
    \begin{equation*}
    y_t = \delta + \phi_1 y_{t-1} + \epsilon_t, \quad \epsilon_t \sim \mathcal{WN}(0, \sigma^2)
    \end{equation*}
    derive the unconditional mean $E[y_t]$ and the unconditional variance $\text{Var}(y_t)$ under the stability condition $|\phi_1| < 1$. What happens if $\phi_1 = 1$?
    
    \item An Augmented Dickey-Fuller test conducted on daily transaction count data yields a test statistic of $-1.82$. If the critical value at the $5\%$ significance level is $-2.89$, state the statistical conclusion. What data transformation is required before fitting an ARMA model?
    
    \item Explain how an adversary executes a ``salami slicing'' fraud scheme (repeatedly siphoning fractions of a cent or a few rupees from millions of customer accounts). Why does standard outlier detection fail to catch this, and how does the CUSUM algorithm detect it?
    
    \item Given an ARIMA$(0, 1, 1)$ model formulated as:
    \begin{equation*}
    y_t - y_{t-1} = \epsilon_t - \theta_1 \epsilon_{t-1}
    \end{equation*}
    show that the one-step-ahead point forecast $\hat{y}_{t+1|t}$ can be expressed as an exponentially weighted moving average (EWMA) of past historical observations.
    
    \item Describe the difference between additive and multiplicative seasonal decomposition in time series. Under what operational conditions in corporate cash flows is multiplicative decomposition required?
\end{enumerate}
